\chapter{Cinemática em duas dimensões}\label{ch:g12:kinematics-2d}

Observe um carro fazendo uma rotatória a constantes \qty{30}{km/h}: o
ponteiro não sai do lugar e, ainda assim, todo passageiro se sente puxado
para o lado --- o movimento está mudando, e não é a \omterm{def:g12:kinematics-2d:velocity}{velocidade escalar}.
Um número por instante já não basta: o movimento no plano precisa de
flechas. Este capítulo equipa um ponto móvel com três --- posição,
velocidade e aceleração --- encadeadas pela derivada, e lê famílias
inteiras de movimentos na geometria delas.

\section{O vetor posição}

O movimento só existe em relação a um \omterm{def:g10:relative-motion:frame}{referencial} declarado --- um corpo
rígido mais um relógio (\cref{ch:g10:relative-motion}). A novidade deste
ano é a precisão: prenda ao \omterm{def:g10:relative-motion:frame}{referencial} uma origem $O$ e dois eixos
graduados perpendiculares $Ox$ e $Oy$, e registre onde o ponto está a
cada instante, como um par de funções.

\begin{definition}[Vetor posição]\label{def:g12:kinematics-2d:position}
Num \omterm{def:g10:relative-motion:frame}{referencial} equipado com uma origem $O$ e eixos $Ox$ e $Oy$, o
\emph{vetor posição}\index{vetor posição} de um ponto móvel $M$ no
instante $t$ é $\vect{OM}(t)$, de coordenadas
$\bigl(x(t),\, y(t)\bigr)$: duas funções do tempo, uma por eixo. A curva
descrita por $M$ enquanto $t$ corre é a sua \omterm{def:g10:relative-motion:trajectory}{trajetória}.
\end{definition}

\section{O vetor velocidade}

\begin{definition}[Vetor velocidade]\label{def:g12:kinematics-2d:velocity}
Ao longo de um intervalo curto o ponto se desloca de
$\vect{OM}(t+\Delta t) - \vect{OM}(t)$; divida por $\Delta t$ e faça o
intervalo encolher: é a derivada do seu curso de matemática, aplicada a
cada coordenada. O \emph{vetor velocidade}\index{vetor velocidade} de $M$
é a derivada do seu \omterm{def:g12:kinematics-2d:position}{vetor posição},
$\vect v(t) = \bigl(x'(t),\, y'(t)\bigr)$, em \unit{m/s}. A norma dele,
$v = \sqrt{x'^2 + y'^2}$, é a \emph{velocidade escalar}\index{velocidade
escalar} --- o único número que o velocímetro mostra.
\end{definition}

\begin{example}[Um drone em voo nivelado]\label{ex:g12:kinematics-2d:drone}
Um drone voa com $x(t) = 4.0\,t$ e $y(t) = 3.0\,t$ (\omterm{def:g10:orders-of-magnitude:unit}{metros}, segundos):
\omterm{def:g10:relative-motion:trajectory}{trajetória} sobre a reta $y = \tfrac34 x$, velocidade
$\vect v = (4.0, 3.0)$ em todos os instantes e \omterm{def:g12:kinematics-2d:velocity}{velocidade escalar}
$\sqrt{4.0^2 + 3.0^2} = \qty{5.0}{m/s}$, constante.
\end{example}

\begin{proposition}[Tangente à trajetória]\label{prop:g12:kinematics-2d:tangent}
A cada instante, $\vect v(t)$ é tangente à \omterm{def:g10:relative-motion:trajectory}{trajetória} na posição atual de
$M$ e aponta no sentido do movimento.
\end{proposition}

\begin{proof}
O deslocamento é uma corda da \omterm{def:g10:relative-motion:trajectory}{trajetória}; à medida que $\Delta t$
encolhe, a direção da corda vira a da tangente.
\end{proof}

\section{O vetor aceleração}

\begin{definition}[Vetor aceleração]\label{def:g12:kinematics-2d:acceleration}
O \emph{vetor aceleração}\index{vetor aceleração} de $M$ é a derivada do
seu \omterm{def:g12:kinematics-2d:velocity}{vetor velocidade},
$\vect a(t) = \bigl(v_x'(t),\, v_y'(t)\bigr) =
\bigl(x''(t),\, y''(t)\bigr)$, em \unit{m/s^2}: com que rapidez, e em que
direção, o \omterm{def:g12:kinematics-2d:velocity}{vetor velocidade} está mudando naquele momento.
\end{definition}

\begin{remark}[Acelerar sem ganhar velocidade]\label{rem:g12:kinematics-2d:turning}
$\vect a \neq \vect 0$ sempre que o \emph{vetor} velocidade muda --- em
comprimento \emph{ou em direção}: um carro freando acelera (para trás);
um carro fazendo curva a \omterm{def:g12:kinematics-2d:velocity}{velocidade escalar} constante acelera (para o
lado). É essa a grandeza que as \omterm{def:g10:inertia:force}{forças} controlam
(\cref{ch:g11:forces-and-motion,ch:g12:newtons-laws}).
\end{remark}

\begin{method}[Ler uma cronofotografia]\label{met:g12:kinematics-2d:chrono}
Uma \emph{cronofotografia}\index{cronofotografia} registra as posições
$M_0, M_1, M_2, \dots$ de um ponto móvel em intervalos de tempo iguais
$\Delta t$.
\begin{enumerate}
  \item O espaçamento dos pontos revela a \omterm{def:g12:kinematics-2d:velocity}{velocidade escalar}:
        espaçamentos iguais, velocidade constante; abrindo, acelerando;
        fechando, freando.
  \item A velocidade em $M_i$ é a corda simétrica,
        $\vect v_i \approx \vect{M_{i-1}M_{i+1}}/(2\,\Delta t)$,
        desenhada em $M_i$, tangente ao caminho pontilhado.
  \item Para a aceleração em $M_i$: copie $\vect v_{i-1}$ e
        $\vect v_{i+1}$ para $M_i$ e desenhe
        $\Delta\vect v = \vect v_{i+1} - \vect v_{i-1}$, de ponta a
        ponta; então $\vect a_i \approx \Delta\vect v/(2\,\Delta t)$.
\end{enumerate}
\end{method}

\begin{omfigure}
\begin{tikzpicture}[scale=1.05, font=\small]
  \fill (0,2.05) circle (2pt) (0.8,1.95) circle (2pt)
    (1.6,1.72) circle (2pt) (2.4,1.36) circle (2pt)
    (3.2,0.87) circle (2pt) (4.0,0.25) circle (2pt);
  \node[above] at (0,2.05) {$M_0$};
  \node[above] at (1.6,1.72) {$M_2$};
  \node[above right] at (2.4,1.36) {$M_3$};
  \node[above right] at (3.2,0.87) {$M_4$};
  \draw[-Stealth, omDef, thick] (1.6,1.72) -- ++(1.2,-0.44)
    node[right] {$\vect v_2$};
  \draw[-Stealth, omDef, thick] (3.2,0.87) -- ++(1.2,-0.83)
    node[right] {$\vect v_4$};
  \draw[-Stealth, omDef, dashed] (2.4,1.36) -- ++(1.2,-0.44);
  \draw[-Stealth, omDef] (2.4,1.36) -- ++(1.2,-0.83);
  \draw[-Stealth, omThm, very thick] (3.6,0.92) -- ++(0,-0.39)
    node[right] {$\Delta\vect v$};
  \draw[-Stealth, omThm, very thick] (2.4,1.36) -- ++(0,-0.75)
    node[left] {$\vect a$};
\end{tikzpicture}

{\small \omterm{met:g12:kinematics-2d:chrono}{Cronofotografia} de uma bola arremessada: $\vect v_2$ e
$\vect v_4$, construídos a partir das cordas, são copiados para $M_3$; a
diferença deles, $\Delta\vect v$, entrega $\vect a$ --- vertical, para
baixo.}
\end{omfigure}

\section{Movimentos retilíneos}

Ao longo de uma reta basta um eixo: $x(t)$, $v(t) = x'(t)$ e
$a(t) = v'(t)$, todos números com sinal.

\begin{definition}[Movimento retilíneo uniforme]\label{def:g12:kinematics-2d:uniform}
Um movimento é \emph{retilíneo uniforme}\index{movimento retilíneo
uniforme} quando $\vect v$ é um vetor constante: \omterm{def:g10:relative-motion:trajectory}{trajetória} reta,
\omterm{def:g12:kinematics-2d:velocity}{velocidade escalar} constante, $\vect a = \vect 0$. Então
$x(t) = v\,t + x_0$.
\end{definition}

\begin{proposition}[Movimento retilíneo uniformemente acelerado]\label{prop:g12:kinematics-2d:uarm}
\index{movimento uniformemente acelerado}%
Se um ponto se move ao longo de $Ox$ com aceleração \emph{constante} $a$,
partindo em $t = 0$ da posição $x_0$ com velocidade $v_0$, então
\[
v(t) = a\,t + v_0,
\qquad
x(t) = \tfrac12\,a\,t^2 + v_0\,t + x_0 .
\]
\end{proposition}

\begin{proof}
$v$ é uma primitiva da constante $a$: $v(t) = a t + C$, e $t = 0$ obriga
$C = v_0$. Por sua vez, $x$ é uma primitiva de $a t + v_0$:
$x(t) = \tfrac12 a t^2 + v_0 t + C'$, com $C' = x_0$.
\end{proof}

\begin{omfigure}
\begin{tikzpicture}[font=\small]
\begin{axis}[omaxis, name=ax1, width=8.4cm, height=3.2cm,
    ylabel={$x$ (\unit{m})}, xmin=0, xmax=4, ymin=0, ymax=21]
  \addplot[omDef, thick, domain=0:4, samples=40] {0.5*2*x^2 + x};
\end{axis}
\begin{axis}[omaxis, name=ax2, at={(ax1.below south west)},
    anchor=north west, yshift=-3mm, width=8.4cm, height=3.2cm,
    ylabel={$v$ (\unit{m/s})}, xmin=0, xmax=4, ymin=0, ymax=10]
  \addplot[omThm, thick, domain=0:4, samples=2] {2*x + 1};
\end{axis}
\begin{axis}[omaxis, at={(ax2.below south west)}, anchor=north west,
    yshift=-3mm, width=8.4cm, height=3.2cm, xlabel={$t$ (\unit{s})},
    ylabel={$a$ (\unit{m/s^2})}, xmin=0, xmax=4, ymin=0, ymax=3]
  \addplot[omProp, thick, domain=0:4, samples=2] {2};
\end{axis}
\end{tikzpicture}

{\small \omterm{prop:g12:kinematics-2d:uarm}{Movimento uniformemente acelerado} ($a = \qty{2.0}{m/s^2}$,
$v_0 = \qty{1.0}{m/s}$, $x_0 = 0$): cada gráfico é a derivada do de cima
--- uma parábola, a inclinação dela e a inclinação da inclinação.}
\end{omfigure}

\begin{example}[Distância de frenagem]\label{ex:g12:kinematics-2d:braking}
Um carro com \omterm{def:g12:kinematics-2d:velocity}{velocidade escalar} $v_0$ freia com desaceleração constante
$a$ (aceleração $-a$). Ele para quando $v(t) = -a t + v_0 = 0$, em
$t_s = v_0/a$, tendo percorrido
$d = -\tfrac12 a t_s^2 + v_0 t_s = v_0^2/(2a)$. O quadrado é a manchete
da segurança viária: dobrar a velocidade \emph{quadruplica} a distância
de frenagem --- com $a = \qty{6.0}{m/s^2}$, são \qty{16}{m} a partir de
\qty{50}{km/h} e \qty{64}{m} a partir de \qty{100}{km/h}.
\end{example}

\section{Movimento circular uniforme e a base de Frenet}

\begin{definition}[Movimento circular uniforme]\label{def:g12:kinematics-2d:ucm}
Um movimento é \emph{circular uniforme}\index{movimento circular
uniforme} quando a \omterm{def:g10:relative-motion:trajectory}{trajetória} é um círculo de raio $R$ e a \omterm{def:g12:kinematics-2d:velocity}{velocidade
escalar} $v$ é constante. O \emph{período}\index{período!de rotação} $T$ é
a duração de uma volta, e a \emph{frequência}\index{frequência!de
rotação} $f = 1/T$ (em \unit{Hz}) é o número de voltas por segundo:
$v = 2\pi R/T = 2\pi R f$.
\end{definition}

\omterm{def:g12:kinematics-2d:velocity}{Velocidade escalar} constante --- e, no entanto, a velocidade gira:
\emph{existe} uma aceleração. Para que lado, e de que tamanho?

\begin{theorem}[Aceleração centrípeta]\label{thm:g12:kinematics-2d:centripetal}
\index{aceleração centrípeta}%
Num \omterm{def:g12:kinematics-2d:ucm}{movimento circular uniforme} de raio $R$ e \omterm{def:g12:kinematics-2d:velocity}{velocidade escalar} $v$, a
aceleração --- chamada de \emph{centrípeta} --- aponta do ponto móvel
para o centro do círculo, com norma constante
\[
a = \frac{v^2}{R}.
\]
\end{theorem}

\begin{proof}
Centre os eixos no centro do círculo e ponha $\omega = v/R$, de modo que
o arco $v t$ subentende o ângulo $\omega t$: a posição é
$x(t) = R\cos(\omega t)$, $y(t) = R\sin(\omega t)$. Derivando,
$\vect v = (-R\omega\sin(\omega t),\, R\omega\cos(\omega t))$, de norma
constante $R\omega = v$; derivando outra vez,
$\vect a = -\omega^2\,\vect{OM}$: oposto ao \omterm{def:g12:kinematics-2d:position}{vetor posição} --- em direção
ao centro --- e de norma $\omega^2 R = v^2/R$.
\end{proof}

\begin{omfigure}
\begin{tikzpicture}[scale=1.0, font=\small]
  \draw[thick] (0,0) circle (1.7);
  \fill (0,0) circle (1.5pt); \node[below left] at (0,0) {$O$};
  \draw[dashed, omExo] (0,0) -- (215:1.7)
    node[midway, below right] {$R$};
  \fill (35:1.7) circle (2pt); \node[above right] at (35:1.7) {$M$};
  \draw[-Stealth, omDef, thick] (35:1.7) -- ++(-0.69,0.98)
    node[above] {$\vect v$};
  \draw[-Stealth, omThm, very thick] (35:1.7) -- ++(215:0.95)
    node[below right] {$\vect a$};
\end{tikzpicture}

{\small \omterm{def:g12:kinematics-2d:ucm}{Movimento circular uniforme}: $\vect v$ tangente e de comprimento
fixo; $\vect a$ apontado para o centro e de comprimento fixo $v^2/R$ ---
girando a velocidade para sempre, sem nunca esticá-la.}
\end{omfigure}

\begin{definition}[Base de Frenet]\label{def:g12:kinematics-2d:frenet}
A \emph{base de Frenet}\index{base de Frenet} na posição atual de um
ponto sobre um caminho curvo é o par de vetores unitários
perpendiculares $(\vect{u_t}, \vect{u_n})$: $\vect{u_t}$ tangente, no
sentido do movimento; $\vect{u_n}$ normal, apontando para o centro do
círculo de raio $R$ que melhor se ajusta à curva ali.
\end{definition}

\begin{proposition}[Aceleração na base de Frenet]\label{prop:g12:kinematics-2d:frenetacc}
Ao longo de qualquer caminho plano de raio local $R$, percorrido com
\omterm{def:g12:kinematics-2d:velocity}{velocidade escalar} $v(t)$,
\[
\vect a = a_t\,\vect{u_t} + a_n\,\vect{u_n},
\qquad a_t = \frac{\dd v}{\dd t},
\qquad a_n = \frac{v^2}{R}:
\]
a \emph{aceleração tangencial}\index{aceleração tangencial} $a_t$ muda a
\omterm{def:g12:kinematics-2d:velocity}{velocidade escalar}; a \emph{aceleração normal}\index{aceleração normal}
$a_n$ gira a velocidade.
\end{proposition}
\admitted

\begin{remark}[Onde mora a demonstração]\label{rem:g12:kinematics-2d:frenetproof}
Demonstramos os extremos: $a_t$ puro numa reta
(\cref{prop:g12:kinematics-2d:uarm}) e $a_n$ puro num círculo
(\cref{thm:g12:kinematics-2d:centripetal}); a decomposição geral é
deduzida no volume do primeiro ano de graduação.
\end{remark}

\begin{example}[Frear entrando na curva]\label{ex:g12:kinematics-2d:bend}
Um carro entra numa curva de raio $R = \qty{50}{m}$ a $v = \qty{15}{m/s}$
enquanto freia a $\dd v/\dd t = \qty{-3.0}{m/s^2}$:
$a_t = \qty{-3.0}{m/s^2}$, $a_n = 15^2/50 = \qty{4.5}{m/s^2}$, de modo
que $a = \sqrt{3.0^2 + 4.5^2} \approx \qty{5.4}{m/s^2}$, apontando para
trás \emph{e} para dentro da curva. Os pneus precisam fornecer as duas
coisas ao mesmo tempo --- e é por isso que os pilotos de corrida freiam
\emph{antes} da curva.
\end{example}

\section{Exercícios}

\begin{exercise}[$\star$]\label{exo:g12:kinematics-2d:1}
Para $x(t) = 3.0\,t$ e $y(t) = 4.0\,t$ (\omterm{def:g10:orders-of-magnitude:unit}{metros}, segundos), dê $\vect v$,
a \omterm{def:g12:kinematics-2d:velocity}{velocidade escalar}, a \omterm{def:g10:relative-motion:trajectory}{trajetória} e a classe do movimento.
\end{exercise}

\begin{exercise}[$\star$]\label{exo:g12:kinematics-2d:2}
Numa \omterm{met:g12:kinematics-2d:chrono}{cronofotografia} feita a cada \qty{0.10}{s}, os pontos estão sobre
uma reta, igualmente espaçados de \qty{0.35}{m}. Identifique o movimento
e calcule a \omterm{def:g12:kinematics-2d:velocity}{velocidade escalar}. O que espaçamentos que se abrem
indicariam?
\end{exercise}

\begin{exercise}[$\star$]\label{exo:g12:kinematics-2d:3}
Para $x(t) = 2.0\,t^2$ ao longo de $Ox$ (\omterm{def:g10:orders-of-magnitude:unit}{metros}, segundos), obtenha
$v(t)$ e $a(t)$, nomeie a classe do movimento e calcule a \omterm{def:g12:kinematics-2d:velocity}{velocidade
escalar} em $t = \qty{3.0}{s}$.
\end{exercise}

\begin{exercise}[$\star$]\label{exo:g12:kinematics-2d:4}
Um cavalinho de carrossel percorre um círculo de raio \qty{4.5}{m}, uma
volta a cada \qty{8.0}{s}: calcule a \omterm{def:g12:kinematics-2d:ucm}{frequência}, a \omterm{def:g12:kinematics-2d:velocity}{velocidade escalar} e a
aceleração (norma e direção).
\end{exercise}

\begin{exercise}[$\star$]\label{exo:g12:kinematics-2d:5}
Um gráfico de $v(t)$ sobe em linha reta de \qty{2.0}{m/s} em $t = 0$ até
\qty{10.0}{m/s} em $t = \qty{4.0}{s}$. Leia a aceleração e depois a
distância percorrida (área sob o gráfico).
\end{exercise}

\begin{exercise}[$\star\star$]\label{exo:g12:kinematics-2d:6}
Um trem de alta velocidade faz uma curva de raio \qty{6.0}{km} a
\qty{320}{km/h}. Calcule $a_n$ e compare-a com $g$; por que as linhas de
alta velocidade precisam evitar curvas fechadas, mesmo perfeitamente
inclinadas?
\end{exercise}

\begin{exercise}[$\star\star$]\label{exo:g12:kinematics-2d:7}
Um carrinho num trilho reto tem $x(t) = -1.2\,t^2 + 8.0\,t + 3.0$
(\omterm{def:g10:orders-of-magnitude:unit}{metros}, segundos). Encontre $v(t)$ e $a$, o instante e a posição da
inversão de sentido, e descreva o movimento antes e depois.
\end{exercise}

\begin{exercise}[$\star\star$]\label{exo:g12:kinematics-2d:8}
Uma \omterm{met:g12:kinematics-2d:chrono}{cronofotografia} ao longo de uma reta, a cada
$\Delta t = \qty{0.20}{s}$, dá $x = 0$, $0.10$, $0.28$, $0.54$ e
$\qty{0.88}{m}$. Estime a velocidade nos três pontos interiores
(\cref{met:g12:kinematics-2d:chrono}) e depois a aceleração; conclua.
\end{exercise}

\begin{exercise}[$\star\star$]\label{exo:g12:kinematics-2d:9}
A Lua orbita a Terra num quase-círculo de raio \qty{3.84e8}{m} em
\qty{27.3}{d}. Calcule a \omterm{def:g12:kinematics-2d:velocity}{velocidade escalar} e a \omterm{thm:g12:kinematics-2d:centripetal}{aceleração centrípeta}
dela; compare a segunda com $g$, sabendo que a Lua está a cerca de $60$
raios terrestres (\cref{ch:g10:universal-gravitation}).
\end{exercise}

\begin{exercise}[$\star\star$]\label{exo:g12:kinematics-2d:10}
O tambor de uma máquina de lavar, de raio \qty{25}{cm}, gira a $1200$
rotações por minuto. Calcule $f$, $T$, a \omterm{def:g12:kinematics-2d:velocity}{velocidade escalar} na borda e a
aceleração na borda em múltiplos de $g$. Por que a roupa fica grudada
enquanto a água sai?
\end{exercise}

\begin{exercise}[$\star\star$]\label{exo:g12:kinematics-2d:11}
Um carro a \qty{90}{km/h} freia com $a = \qty{7.5}{m/s^2}$ constante:
calcule o tempo de parada e a distância de frenagem, e esboce $v(t)$.
\end{exercise}

\begin{exercise}[$\star\star\star$]\label{exo:g12:kinematics-2d:12}
Para $x(t) = R\cos(\omega t)$ e $y(t) = R\sin(\omega t)$: calcule
$\vect v$; verifique que a \omterm{def:g12:kinematics-2d:velocity}{velocidade escalar} é a constante $R\omega$ e
que $\vect v \perp \vect{OM}$; calcule $\vect a$, mostre que ele vale
$-\omega^2\,\vect{OM}$ e recupere $a = v^2/R$.
\end{exercise}

\begin{exercise}[$\star\star\star$]\label{exo:g12:kinematics-2d:13}
Uma motocicleta entra numa curva de raio \qty{80}{m} a \qty{20}{m/s},
freando a \qty{2.5}{m/s^2}. Calcule $a_t$, $a_n$, a norma de $\vect a$ e
o ângulo dele com a direção do movimento. Os pneus aderem até
\qty{7.0}{m/s^2} no total: o piloto está dentro do limite?
\end{exercise}

\begin{exercise}[$\star\star\star$]\label{exo:g12:kinematics-2d:14}
Um metrô acelera do repouso a \qty{1.2}{m/s^2} até \qty{20}{m/s}, segue
em velocidade de cruzeiro por \qty{40}{s} e depois freia a
\qty{1.5}{m/s^2} até parar. Calcule a duração e o comprimento de cada
fase, a distância total e a \omterm{def:g10:relative-motion:speed}{velocidade média}. Esboce $v(t)$ e confira a
distância como uma área.
\end{exercise}

\begin{exercise}[$\star\star\star$]\label{exo:g12:kinematics-2d:15}
Um \omterm{def:g10:universal-gravitation:satellite}{satélite} geoestacionário gira num raio de \qty{4.22e7}{m} a partir do
centro da Terra, com \omterm{def:g12:kinematics-2d:ucm}{período} \qty{86164}{s}. Calcule a \omterm{def:g12:kinematics-2d:velocity}{velocidade escalar}
e a \omterm{thm:g12:kinematics-2d:centripetal}{aceleração centrípeta} dele. O que a fornece, e por que o \omterm{def:g10:universal-gravitation:satellite}{satélite} fica
pairando sobre um ponto fixo do equador
(\cref{ch:g12:satellites-kepler})?
\end{exercise}

\section{Problema: o laudo do acidente}

\begin{problem}[{Problema de fim de semana --- o laudo do acidente: um
vídeo de câmera de bordo lido quadro a quadro, uma lei de frenagem
reconstruída a partir das marcas de derrapagem, uma rotatória que
testemunha sobre a aderência e um veredito em quilômetros por hora}]
\label{pb:g12:kinematics-2d:1}
Um carro atinge uma van que saiu de uma rua \omterm{def:g12:mechanical-waves:transverse}{transversal}; ninguém se
machuca, mas as seguradoras discordam. O perito tem imagens de uma câmera
aérea (um quadro a cada \qty{0.50}{s}), marcas de derrapagem medidas e a
câmera de bordo do carro. O limite de velocidade é \qty{50}{km/h}.

\medskip
\noindent\textbf{Parte I --- Ler o filme.}
Os quadros anteriores a qualquer reação dão, ao longo da estrada reta:
\begin{center}\small
\begin{tabular}{lccccc}
$t$ (\unit{s}) & 0.00 & 0.50 & 1.00 & 1.50 & 2.00\\
$x$ (\unit{m}) & 0.0 & 9.0 & 18.0 & 27.0 & 36.0
\end{tabular}
\end{center}
\begin{enumerate}
  \item Por que o laudo precisa nomear primeiro um \omterm{def:g10:relative-motion:frame}{referencial} e eixos?
        Quais são usados aqui?
  \item Calcule a \omterm{def:g10:relative-motion:speed}{velocidade média} em cada um dos quatro intervalos.
  \item Que classe de movimento é essa? Justifique pela tabela.
  \item Dê a velocidade $v_0$ do carro em \unit{m/s} e em \unit{km/h}.
  \item Qual é a aceleração ao longo desses dois segundos?
\end{enumerate}

\medskip
\noindent\textbf{Parte II --- As marcas de derrapagem.}
Testes de pneus nesse asfalto seco dão uma desaceleração de frenagem
constante $a = \qty{6.0}{m/s^2}$; as marcas de derrapagem têm
$d_b = \qty{27}{m}$ de comprimento. Tome $t = 0$ e $x = 0$ no início da
frenagem, com velocidade desconhecida $v_B$.
\begin{enumerate}[resume]
  \item Por primitivas, escreva $v(t)$ e $x(t)$ durante a frenagem.
  \item Exprima o tempo de parada $t_s$; mostre que $d_b = v_B^2/(2a)$.
  \item Deduza $v_B$ a partir das marcas de derrapagem. É coerente com a
        Parte I?
  \item O tempo de reação é \qty{1.0}{s}: calcule a distância de reação e
        a distância total de parada.
  \item Explique, a partir da questão 7, por que dobrar a velocidade
        quadruplica a distância de frenagem; calcule $d_b$ a exatamente
        \qty{50}{km/h}.
\end{enumerate}

\medskip
\noindent\textbf{Parte III --- A rotatória testemunha.}
Um pouco antes, a câmera de bordo mostra o carro percorrendo uma faixa de
rotatória de raio $R = \qty{12.5}{m}$ com \omterm{def:g12:kinematics-2d:velocity}{velocidade escalar} constante,
um quarto de volta em \qty{2.5}{s}. A pista consegue fornecer no máximo
$a_n = \mu g$, com $\mu = 0.70$.
\begin{enumerate}[resume]
  \item A \omterm{def:g12:kinematics-2d:velocity}{velocidade escalar} ali é constante --- e a velocidade? O carro
        está acelerando? Explique.
  \item Calcule a \omterm{def:g12:kinematics-2d:velocity}{velocidade escalar} do carro na rotatória.
  \item Calcule a \omterm{thm:g12:kinematics-2d:centripetal}{aceleração centrípeta} dele (norma e direção).
  \item Calcule a velocidade máxima sem derrapagem. O carro estava dentro
        dela?
  \item Dê o \omterm{def:g12:kinematics-2d:ucm}{período} e a \omterm{def:g12:kinematics-2d:ucm}{frequência} de uma volta completa nessa
        velocidade.
\end{enumerate}

\medskip
\noindent\textbf{Parte IV --- O veredito.}
A van apareceu $D = \qty{41}{m}$ à frente do carro no instante em que o
motorista começou a reagir.
\begin{enumerate}[resume]
  \item De quanto, em \unit{km/h} e em porcentagem, a velocidade antes da
        frenagem excedeu o limite?
  \item Calcule a velocidade do carro no impacto.
  \item Refaça a conta de parada a exatamente \qty{50}{km/h} (mesma
        reação, mesma frenagem): o carro para, e com que margem?
  \item Dê a fórmula geral da distância de parada $d(v)$; por que a
        velocidade entra duas vezes --- uma linearmente e outra ao
        quadrado?
  \item O veredito, em uma frase com os números: a que velocidade o
        motorista estava, e o acidente teria acontecido no limite?
\end{enumerate}
\end{problem}
